\documentclass{article}
\usepackage{amsmath, amssymb, amsthm}
\usepackage{hyperref}
\usepackage{geometry}
\geometry{margin=1in}

\title{Erd\H{o}s Problem \#233}
\date{Last updated: 2025-08-31}
\author{Source: erdosproblems.com}

\begin{document}
\maketitle

\noindent\textbf{Status:} OPEN \quad \textbf{No prize}

\noindent\textbf{Tags:} number theory, primes

\medskip
\noindent\textbf{Problem Statement:}

Let $d_n=p_{n+1}-p_n$, where $p_n$ is the $n$th prime. Prove that\[\sum_{1\leq n\leq N}d_n^2 \ll N(\log N)^2.\]

\bigskip
\noindent\textbf{Remarks:}

Cramer \cite{Cr36} proved an upper bound of $O(N(\log N)^4)$ conditional on the Riemann hypothesis. Selberg \cite{Se43} improved this slightly (still assuming the Riemann hypothesis) to\[\sum_{1\leq n\leq N}\frac{d_n^2}{n}\ll (\log N)^4.\]The prime number theorem immediately implies a lower bound of\[\sum_{1\leq n\leq N}d_n^2\gg N(\log N)^2.\]This would imply in particular that $d_n\ll n^{1/2}\log n$ for all $n$, which is known only the assumption of the Riemann Hypothesis.

The values of the sum are listed at A074741 on the OEIS.

This is discussed in problem A8 of Guy's collection \cite{Gu04}.

\bigskip
\noindent\textbf{References:}

[Cr36] Cram\'{e}r, Harald, On the order of magnitude of the difference between consecutive prime numbers. Acta Arithmetica (1936), 23--46.
 
 [Gu04] Guy, Richard K., Unsolved problems in number theory. (2004), xviii+437.
 
 [Se43] Selberg, Atle, On the normal density of primes in small intervals, and the
difference between consecutive primes. Arch. Math. Naturvid. (1943), 87--105.

\bigskip
\noindent\small{Source: \url{https://www.erdosproblems.com/233}}
\end{document}
